\paragraph{\textbf{General framework of supervised ML}}
In the most general from, supervised ML models $\bm{y}=f_{\boldsymbol{\theta}}(\bm{x})$ are constructed to achieve two types of tasks, \textbf{\textit{regression}} and \textbf{\textit{classification}}, which depends on the type of the target variable $\bm{y}$. Generally, $\forall y_i\in\bm{y}$, $y_i\in\mathbb{R}$ takes the form of a continuous variable in regression, while $y_i\in\mathbb{Z}^{+}$ takes discrete values in the classification task, where each class is represented by a positive integer. Without loss of generality, we will focus on regression in the following discussion, and the same principles can be applied to classification by simply changing the target distributions from a continuous one into one for discrete random variables [see \cref{appendix:logistic_regression}].

At the core of the design of ML/AI models is how to construct the loss function $\mathcal{L}(\boldsymbol{\theta})$ that is used to train the model. Owing to the simplicity of the supervised ML framework, that is fundamentally identical to linear regressions, most works on training the supervised ML models adapted a $L2$ loss function without in-depth discussions, which minimises the mean squared error (MSE) between the model predictions and the target values to obtain a single set of optimum parameters (that is also widely used in training traditional atom-atom forcefields, see \cref{sect:empirical_forcefield}):
\begin{equation}
    \label{eq:simple_ML_optimisation}
    \hat{\boldsymbol{\theta}}=\argmin_{\boldsymbol{\theta}}\mathcal{L}(\boldsymbol{\theta})=\argmin_{\boldsymbol{\theta}}\left\Vert \bm {y}-f_{\boldsymbol{\theta}}(\bm{x})\right\Vert^2_2=\argmin_{\boldsymbol{\theta}}\sum_{i=1}^N (y_i-f_{\boldsymbol{\theta}}(x_i))^2.
\end{equation}


\begin{figure}[tb]
    \centering
    \includegraphics[width=\linewidth]{ridge_regression_prob_dist_loss_landscape.pdf}
    \caption{Probabilistic view of hyperparameter optimisation in rigde regression. (a) The loss function landscape as the function of the two hyperparameters $\phi_1$ and $\phi_2$, where the high (low) loss values are shown in yellow (green). (b) The global minimum of the loss function, which on the left panel shows the ground truth (dashed line) and predicted (solid line) Gabor function $y=f(x|\phi_0,\phi_1)$ perfectly overlapping with each other, and the training data (yellow dots) that are sampled from the ground truth function with noise. The right panel shows the probability distribution of the hyperparameters $\phi_0$ and $\phi_1$, which, at the global minimum, shows a narrow distribution with almost zero off-diagonal covariance. Panels (c-f) show the same information as (b) but computed at different points in the hyperparameter space, which are marked correspondingly in (a). To facilitate cross comparison, all probability distributions are plotted with $x$ and $y$ axis in the range of $[\hat{\phi}-1,\hat{\phi}+1]$ where $\hat{\phi}$ is the medium of the normal distribution. It can be seen that the probability distributions of the hyperparameters become broader and strongly correlated at the points where the loss function has high values. For details, see \cref{appendix:details_ridge_regression_loss_landscape}. [The underlying example is taken from Fig. 6.4 in \citet{prince2023understanding}]} 
    \label{fig:ridge_regression_prob_dist_loss_landscape}
\end{figure} 

However, as we have thoroughly discussed with more details in \cref{appendix:linear_regression}, the loss function can be derived from a probabilistic perspective. This perspective is established based on two equivalent viewpoints for training the supervised ML models, whereby model training can be viewed as a process to reach a probability distribution: (i) of the observables that maximises the likelihood for observing the training observables conditioned on corresponding model inputs and parameters [see \cref{eq:bayes_linear_variational_inf_expectation_variational_dist}, which in \cref{sec:ML_AI_PES_sampling_info_theory} we have also discussed extensively how this is related to statistical samplings on PES.] or, (ii) of the model hyperparameters to one that maximises the evidence for observing the training data conditioned on the model inputs [see \cref{eq:bayes_linear_variational_inf}, in which the integrand in the marginal distribution can be essentially bought down to a delta function $\delta(\boldsymbol{\theta}-\hat{\boldsymbol{\theta}})$ or alternatively a narrow Gaussian distribution around $\hat{\boldsymbol{\theta}}$]. The concept behind the second perspective is further illustrated in \cref{fig:ridge_regression_prob_dist_loss_landscape} by showcasing how the probability distributions of the model hyperparameters evolve across different parts of the landscape of the training loss function. 

As elaborated in \cref{appendix:ridge_regression}, one of the key mathematical tools for us to establish the second viewpoint for deriving the concrete form of the loss function to enable model training is the utilisation of \emph{variational Bayesian inference}, which is a powerful framework for approximating complex probability distributions through the introduction of a variational distribution $q(\boldsymbol{\theta})$ to approximate the true posterior distribution $p(\boldsymbol{\theta}|\bm{y},\bm{x})$ of the model parameters. This variational distribution also provides us a systematic way to encode our prior beliefs about the model parameters, which can be crucial for regularising the model training process to enhance the model stabilities and generalisation capabilities through minimising the tendency of overfitting the training data, particularly the outliers. Introducing a distribution over the model parameters also allows us to capture the uncertainty in the model predictions, which is essential for making informed decisions based on the model outputs. This can be either achieved in linear ridge regression [\cref{appendix:ridge_regression}], or more complex nonlinear models such as the Kernel Ridge Regression (KRR) [\cref{appendix:krr}] and the Gaussian process  [\cref{appendix:GP}] models. Finally, the variational Bayesian inference framework, particualy, the training process through maximising the evidence for observing the training data $p(\bm{y}|\bm{X})$ that can be casted into a problem of maximising the evidence lower bound (ELBO) [see \cref{eq:bayes_linear_variational_inf_ELBO}], provides us with the unified foundations for understanding the training process of a wide range of ML/AI models,\cite{prince2023understanding} particularly for the generative models, such as variational autoencoders (VAEs)\cite{diederik2019introduction} and diffusion models.\cite{sohl2015deep,song2019generative} For this reason, we believe that introducing the variational Bayesian inference framework to the readers within the context of simple linear regression is ideal for building the foundations for understanding the training process of more complex ML/AI models.

Based on the above conceptual framework, without diverting the readers into in-depth mathematical details, we provided readers who are interested these details in the \cref{appendix:linear_regression} on how to derive the loss function from the perspective of variational Bayesian inference for a few promoinent linear regression models in the context of materials and molecular modelling, and provides some general highlights on the key features of the models that have been discussed in \cref{appendix:linear_regression}. 

We begin from the simplest linear regression model, which is the ordinary least squares (OLS) regression, that aims to fit a linear relationship $y=\bm{x}^\top\boldsymbol{\beta}$ \footnote{Following the convention in statistics, we use $\boldsymbol{\beta}$ to denote the model parameters for linear regression, which is equivalent to $\boldsymbol{\theta}$ used in the general framework of supervised ML.} between the input features $\bm{x}$ and the target variable $y$, where by modelling both the target and predicted outcomes as random variables that follows Gaussian distributions, we cast the model optimisation problem into one that minimises the KL divergence between these two distributions, which naturally leads to the analytical form shown in \cref{eq:simple_ML_optimisation}. 

To avoid the model from overly weighting a few components $\beta_i\in\boldsymbol{\beta}$, which can be particularly problematic when the input features are highly correlated, ridge regression [\cref{appendix:ridge_regression}] introduces an additional regularisation term $\lambda\left\Vert \boldsymbol{\beta}\right\Vert^2_2$ to the loss function to stabilise the model training, where $\lambda$ is a hyperparameter that controls the strength of the regularisation. As shown in \cref{appendix:ridge_regression}, by adapting the variational Bayesian inference framework, we can formally establish the ridge regression process as one that tries to stabilise the OLS through encoding our prior belief on the probability distributions $q(\boldsymbol{\beta})$ that the model parameter should follow into the model training, such that the optimisation will focus on finding the optimal parameters in the region close to $\boldsymbol{\beta}$ provided by our prior belief. 

With the loss function for ridge regression derived based on this fundamental framework [\cref{eq:bayes_linear_variational_inf_loss}], it is easy to modify it to the corresponding loss function for (i) a classification task, such as logistic regression [\cref{appendix:logistic_regression}], by changing the conditional distribution of the target variable $p(y|\bm{x},\boldsymbol{\beta})$ from a Gaussian distribution to one that supports discrete variables, such as the Logistic distribution, and (ii) to adopt other forms of regularisation, such as the LASSO  regression, which as shown in \cref{appendix:lasso_regression}, is fundamentally driven by the inspiration of changing the prior distribution of the model parameters from a Gaussian distribution to a Laplace distribution, which has a sharper peak at zero and heavier tails than the Gaussian distribution, thus encouraging sparsity in the model parameters and enhancing the model interpretability by effectively performing feature selection with this type of regularisation. 

To keep the discussion general, in \cref{appendix:multimodal_bayesian_regression}, we also provided a brief account how this framework can be generalised to the case where the model parameters are expected to follow a complex multimodal distribution, such as one that is described by the Gaussian Mixture Model (GMM). Conceptually, we can draw some similarity with the training of classical atom-atom forcefields [\cref{sect:empirical_forcefield}] in which the forcefield parameters are simultaneously optimised against a set of heterogeneous targets $\bm{y}=(E,\mathbf{F},\boldsymbol{\Phi})^\top$ that comprised of energy, atomic forces and the force constant matrix.\cite{grimme2014general,rohskopf2017empirical} In this case, the model output can be described by a mixture of Gaussian distributions each for a different target: 
\begin{equation*}
p(\bm{y}\vert \bm{X},\sigma_E^2,\sigma_{\bm{F}}^2,\sigma_{\boldsymbol{\Phi}}^2)\sim\mathcal{N}\left(E|\bm{X}^\top\boldsymbol{\beta}_E,\sigma_E^2\bm{I}\right)\mathcal{N}\left(\bm{F}|\bm{X}^\top\boldsymbol{\beta}_{\bm{F}},\sigma_\bm{F}^2\bm{I}\right)\mathcal{N}\left(\bm{F}|\bm{X}^\top\boldsymbol{\beta}_{\boldsymbol{\Phi}},\sigma_{\boldsymbol{\Phi}}^2\bm{I}\right),
\end{equation*}
which, upon taking the log-likelihood, automatically leads to a loss function that is a weighted sum of the mean-squared-error for each target for model training:
\begin{equation*}
    \frac{1}{\sigma_E^2}\left\Vert E-\bm{X}^\top\boldsymbol{\beta}_E\right\Vert^2_2+\frac{1}{\sigma_\bm{F}^2}\left\Vert \bm{F}-\bm{X}^\top\boldsymbol{\beta}_{\bm{F}}\right\Vert^2_2+\frac{1}{\sigma_{\boldsymbol{\Phi}}^2}\left\Vert \boldsymbol{\Phi}-\bm{X}^\top\boldsymbol{\beta}_{\boldsymbol{\Phi}}\right\Vert^2_2.
\end{equation*}
However, we are not aware of the \emph{direct} wide practical applications of this framework in which the \emph{model parameters are directly drawn and optimised from a multivariant Gaussian distribution}, particularly in the context of materials and molecular modelling, because of the high computational cost associated with sampling from complex multimodal distributions that is required for the model training. 

To enhance the capabilities to handle nonlinear relationships between the input features and the target variable, the linear regression models can be further extended to their nonlinear variants, which are discussed in \cref{appendix:nonlinear_regression}. More specifically, we begin with the Kernel Ridge Regression (KRR) model [\cref{appendix:krr}] which extends the simple ridge regression by capturing the nonlinear relationships between the input features and the target variable through the utilisation of kernel functions, that implicitly map the input features into a higher-dimensional feature space and accounts for the nonlinear relationships in the feature space through the utilisation of nonlinear kernel functions. The so-called similarity kernel is a central concept in materials and molecular modelling in ML/AI, which we shall further elaborated in later sections.  To avoid the model being biased by a few outliers, the support vector machine (SVM) model [\cref{appendix:SVM}] stabilises the KRR model by introducing additional filters in the loss function to remove outlying points during the model training. Finally, in \cref{appendix:GP}, we provided detailed discussion on the Gaussian Process (GP) regression model from the perspective of variational Bayesian inference that optimises the probability distributions over the model function space rather than the model parameter space, which provides the greatest flexibilities for capturing complex nonlinear relationships between the input features and the target variable, and is one of the most popular supervised ML models in the atomistic modelling community.\cite{deringer2021gaussian}

\paragraph{\textbf{Supervised ML for atomistic materials and molecular modelling}}

Here, we shall briefly review some of the notable applications of supervised ML in atomistic materials modelling that are particularly related to the sampling of PES and physical properties of materials that have been discussed in the previous sections. 

\begin{figure}[tb]
    \centering
    \includegraphics[width=\textwidth]{GAP_forcefield.png}
    \caption{(a) Comparison of the performances for modelling the energy landscapes of a migrating vacancy in GaN (top) and phase transformation of rhombohedral graphite to diamond (bottom) as computed from DFT, the GAP model, and the empirical Tersoff and Brenner forcefields [adopted from \citet{bartok2010gaussian}]. (b) Time-dependent evolution of energy for MAPbI$_3$ perovskite computed with on-the-fly MLFF trained from the underlying DFT model using GP regression. [adopted from \citet{jinnouchi2019phase}], showing that the (re)-training of the MLFF can be triggered when the model uncertainty exceeds a certain threshold.}
    \label{fig:GAP_forcefield}
\end{figure}

\textbf{\textit{Energy Models}} One of the primary goal for applying supervised ML in atomistic modelling is to construct machine-learning force fields (MLFF) that can expedite the statistical samplings of PES with computational costs that is comparable to empirical forcefields, but with the accuracy that is in par with the first-principles calculations. \citet{bartok2010gaussian} pioneered the development of the Gaussian approximation potential (GAP) model, that is combined with the bispectrum form of the density-based representation atomic environment.\cite{bartok2013representing,willatt2019atom} As shown in \cref{fig:GAP_forcefield}(a), the GAP model can accurately capture the variation in the system energy across the defect migration and phase transition processes, which typically involve bond breaking and formation, which failed to be captured by the empirical Tersoff and Brenner forcefields, and would traditionally require the application of more advanced reactive forcefields. This provides a good example to showcase the capabilities of GP as a single framework to capture complex nonlinear relationships across regions of the PES that are fundamentally described by different physics.  \citet{thompson2015spectral} adopted  the spectral neighbour analysis potential (SNAP) model, to learn a linear relationship between energy, forces and stress tensor with the bispectrum representations of the local atomic environments to circumvent the speed bottleneck of nonlinear GP process to accelerate the modelling of large molecules. \citet{glielmo2017accurate} developed covariant kernel that respects the SO(3) rotational equvariance of the atomic forces to enable the direct learning of atomic forces using GP from the underlying DFT data. Finally, \citet{jinnouchi2019phase} developed an on-the-fly PES sampling scheme that is based on the GP regression model, in which the training of the MLFF is performed simultaneously with the sampling process, and the retraining of the MLFF can be triggered when the model uncertainty exceeds a certain threshold, which allows the model to achieve a good balance between accuracy and computational cost during the sampling process [see \cref{fig:GAP_forcefield}(b)]. The incorporation of this approach into the mainstream VASP code\cite{Kresse1996CMS,Kresse1996PRB,Kresse1999PAW} represents a significant milestone for the applications of MLFF in atomistic modelling.

In addition to the application of GP in developing MLFF, many other supervised ML models have also been explored for this purpose. For example, \citet{rupp2012fast} combined the KRR model with the Coulomb matrix representation of the molecular structure to develop a MLFF for small organic molecules, which can be used to predict the atomisation energies of these molecules with high accuracy. This represents one of the earliest applications of supervised ML for predicting the properties of molecules. Fundamentally, this simple model works because the Coulomb matrix representations, albeit its simple mathematical form, can capture the nonlinear atom-atom Coulomb interactions that are the dominant contribution to the total energy of the system, which provided a strong physical prior that is further enhanced by the capabilities to handle nonlinear relationships in the KRR model to achieve accurate predictions. \citet{xie2023ultra} went back to the drawing board to start from the many-body expansion of the total energy [\cref{eq:total_energy_atom_pos_expand}], which uses the B-spline basis functions to construct the feature vectors that captures the 2- and 3-body atom-atom interactions to be fed into the KRR model to develop a MLFF for small organic molecules. \citet{podryabinkin2017active} developed an active learning scheme to train a moment tensor potential (MTP) with linear regressions. Finally, as discussed in \cref{sec:anharmonic_vibs}, compressive sensing has been applied to develop TDEP\cite{hellman2013temperature,zhou2019compressive} for modelling the anharmonic vibrational behaviours of solids at finite temperatures. This is essentially developing a forcefield in which the energy is polynomially expanded with respect to the atomic displacements, and the ML model is trained to learn the coefficients of this expansion, which are the high-order force constant matricies that are used to capture the anharmonicity in the vibrational behaviours of solids.

For training the MLFF, the input $\bm{x}$ is typically some feature vector $\mathcal{X}$ that is constructed solely based on the atomic structure $\bm{R}$. As will be demonstrated in the subsequent discussions, by including other elemental physical properties into the feature vector that is to be combined with supervised ML approach will enables us to predict a wider range of physical properties. %When solving the classical physics problems, such as the establishment of Newton's second law for kinematics and the Ohm's law for electric circuits, the underlying linear relationships among the physical quantities can be trivially revealed from the experimental data. However, as the physical properties that we are interested in become more complex, in which a range of interdependent physical properties may contribute to the target properties collectively, such as the interplay between the electronic and phonon conductivities to the thermoelectric performance of materials, supervised ML models provided us a powerful tools to uncover the complex nonlinear relationships between the input features and the target properties, which can be crucial for facilitating us in understanding the complex relationships among the interdependent physical properties.
Instead of the actual ML model being employed, the focus of the following discussion will be on the key physical information that is used to construct the descriptor input $\mathcal{X}$, that are believed to provide dominant contribution to the target properties being predicted, which are the key enabler that allows complex physical properties to be learned with simple ML models. The rationale behind this approach can be understood by drawing parallel comparison to how classical physics laws were established. For examples, the Newton's second law for kinematics and Ohm's law for electrical circuits, in which there exists a small number of interdependent variables whereby the relationships among them were trivial to be discovered from a small set of experimental data. As the materials' properties that we are interested in become more complex, which are cooperatively determined by a large set of possible factors, some algorithmic approaches, such as ML, becomes indispensable in facilitating up to uncover the complex relationships amongst the physical and chemical properties of materials.

\textbf{\textit{Phononic Properties}} \citet{zhuo2018identifying} applied SVM to predicted the Debye temperatures for phosphors, which is a key proxy properties that can be correlated with the photoluminescence quantum yield of the phosphors, in which the input features are 150 dimensional that comprises a ranges of physical properties, such as the  radii and electronic configurations for the consistituent elements, as well as space group informations, which collectively determines the bonding strengths and the vibrational characteristics of the solids. \citet{juneja2019coupling} incorporated fundamental phonon properties of solids, such as the maximum harmonic phonon frequency, integrated Gr\"{u}neisen parameters (up to 3 THz), that are combined with other elemental and structural information, such as the atomic masses and unit cell volumes, that are subsequently fed into a GP model to accelerate the prediction of lattice thermal conductivities ($\kappa_L$) of solids. This provides an example on how ML models can be used to overcome the sampling challenges to compute the anharmonic phonon properties of solids from first principles that contribute towards $\kappa_L$ [\cref{sec:anharmonic_vibs}]. Similar approaches have been adapted in other works to expedite the determination of $\kappa_L$ for inorganic crystals.\cite{chen2019machine,jaafreh2021lattice}

\textbf{\textit{Electronic Properties}}  \citet{lee2016prediction} applied a set of supervised ML models, including ordinary least square, LASSO and SVR, to predict the $G_0 W_0$ band gaps for 270 inorganic compounds from KS band gaps, cohesive energy, crystalline volume per atom, and other fundamental information of constituent elements, achiving an accuracy of 0.24 eV in RMSE. \citet{zhuo2018predicting} applied a SVM classifier to first separate metals from non-metals for inorganic solids, then used a SVR algorithm to predict the band gaps for the non-metallic solids, in which an 136-dimensional feature vector that was constructed purely from the elemental properties was used as the input for the ML models. \citet{pilania2017multi} presented a multi-fidelity co-kriging statistical learning framework to predict the high-fidelity hybrid-DFT electronic band gap from  the low-fidelity GGA band gaps for 600 elpasolite compounds, in which the co-kriging model is essentially a GP regression model that is trained to learn the correlation between the low-fidelity and high-fidelity data to achieve accurate predictions for the high-fidelity data with reduced computational cost, based on features that include the chemical identities at the A, B, B$'$ and X sites, the electronegativity, first ionisation potentials, emprical atomic raidii and the Pettifor's Mendeleev numbers for the constituent elements. \citet{li2022predicting} applied similar approach to construct feature vectors that are combined with linear and kernel ridge regressions to predict the variations in the band gap energies for ABO$_3$ perovskites that were computed from PBEsol and HSE06 density functionals. Importantly, the authors applied the gradient boosted decision tree (GBDT) model to analyse the importance of different features for predicting the band gap variations, which is particularly useful for enhancing the interpretability of the model, thus facilitating the understanding of the important physical properties that contribute to the band gap variations in these perovskites. Finally, \citet{weston2018machine} developed SVR models to predict the band gap energies for 1568 kesterite compounds for photovoltaic applications using features that were constructed with the SVM model. 

\textbf{\textit{Electron-Phonon Interactions and Superconductivities}} In \cref{sec:eph_coupling}, we have discussed the fundamentals of EPI that plays a critical role in determining the temperature-dependent electronic and superconduncting properties of solids, in which the challenge behind the determinations of the EPI strengths in solids is, in the language of probability theory, the sampling of joint probability distributions of the electronic and phononic densities of states under the mutual interactions of the two subsystems. Hence, the purpose of developing simple ML surrogate models is to directly overcoming this sampling challenge to facilitate the direct learning of EPI properties from more elementary physical properties of materials. For examples, \citet{haldar2024machine} tried to predict the zero-point renormalisation of the electronic bandgaps of 2D semiconductors due to EPI using the Extra Random Tree model that was applied on 9 different features, including the zero-Kelvin electronic bandgap energies, the dielectric constants, as well as fundamental phononic and atomic properties. For similar purpose, \citet{nandi2024machine} applied GP model,  employing more than 37 features that chracterise the physical properties of materials to predict the parameters that is used in the O'Donnell model to predict the temperature-dependent bandgap renormalisation of metal oxides. \citet{tran2023machine} generated materials features using MATMINER\cite{ward2018matminer} to train a GP model to predict the EPI strengths and $\omega_{\mathrm{log}}$  to facilitate the discoveries of conventional superconductors based the McMillan model [c.f. \cref{eq:mcmillan_Tc}]. An interesting approach was put forward by \citet{xie2019functional}, which applied data-driven method using the SISSO\cite{ouyang2018sisso} algorithm to discover alternative functional form to the McMillan model for predicting the superconducting transition temperatures from the fundamental parameters of $\lambda$, $\mu^{*}$ and $\omega_{\mathrm{log}}$. Finally, EPI also plays an important role in attenuating the electronic transport properties of thermoelectric materials. For this, \citet{wee2019quantification} applied GP models to predict the EPI strengths $\left\vert g_{nm\nu}(\bm{k},\bm{q}) \right\vert^2\sim g_\nu^2(\varepsilon_{n\bm{k}},\varepsilon_{m\bm{k}+\bm{q}})$ as a function of the electronic energies $\varepsilon_{n\bm{k}}$ and $\varepsilon_{m\bm{k}+\bm{q}}$ for a range of materials.

\textbf{\textit{Electronic Charge Densities}} Charge density is a key fundamental property of materials and molecules, whereby in the DFT framework, having the access to the charge density is the key for solving the KS equation, from which other auxiliary properties can be computed. In this regards, the ability of predicting charge density using ML approach has fundamental interests for atomistic modellings. Using the symmetry-adapted GP regression (SA-GPR), 
\citet{grisafi2018transferable}  predicted the expansion coefficients to reconstruct the molecular charge densities from atom-centred basis functions, using the SOAP-type\cite{bartok2013representing} structure descriptors to represents the covalent bonding environments in molecules. 
The same approach was further extended to predict the charge densities in non-covalently bounded molecular systems,\cite{fabrizio2019electron} solids,\cite{lewis2021learning} that can be ultimately applied to accelerate the QM/MM simulations of electrochemical interfaces.\cite{grisafi2024accelerating} Ultimately, the goal behind the development of ML model is to try bypassing the need of solving the KS eigen-equations. 
In this regard, \citet{brockherde2017bypassing} developed a KRR model to directly learn the Hohenberg-Kohn map from the external potential to the ground state charge density, which can be used to predict the energies and atomic forces for MD simulations of molecules.
 
\textbf{\textit{Magnetic Properties}} By applying supervised ML models to directly predict the spin configurations and associated magnetic properties of materials can be considered as a promising route to overcoming the high computational costs associated with sampling the spin configurations in solids to determine their likely spin ground states, and derived associated magnetic properties therefrom [\cref{sec:mag_prop}]. 
For instances, \citet{dam2018important} employed the KRR model to predict the Curie temperatures of binary transition metal alloys, utilising elemental descriptors including information such as the $3d$ and $4f$ orbital and spin magnetic moments which play a significant role towards determining the ground state magnetic state of the binary alloys.
\citet{long2021accelerating} combined the ground-state magnetic configurations, chemical and structural information (via the SOAP\cite{bartok2013representing} descriptor), that were fed into the Random Forest model to predict the ground state magnetic configuration and Curie temperatures to accelerate the classifications of ferromagnetic and antiferromagnetic compounds.
\citet{ghosh2020machine} also employed the Random Forest model to predict the spin and orbital magnetic moments of uranium-based compounds. The descriptors that were used for their model learning include information such as structural information, cation $5f$-subshell occupation, which were further augmented with descriptors generated from these primitive information through mapping from 10 prototypical mathematical functions.
Lastly, the Random Forest model was also constructed in the study by \citet{khatri2023magnetization} to predict the magnetic moment per ion in Fe-based compounds, in which the orbital field matrix was used for representing the structures of the compounds. To enhance the model interpretability, the authors also performed SHAP (SHapley Additive exPlanations) feature importance analysis\cite{lundberg2017unified} to understand the contributions of orbital interactions in the feature vectors towards the accuracies of model predictions.

\textbf{\textit{Polarisations, Ferro- and Piezo-electric Properties}} \citet{xue2016accelerated} presented one of the earliest works in this area to help the design of \ce{BaTiO3}-based solid solutions as piezoelectric materials by constructing the phase diagram based on the Landau-Devonshire theory. In particular, Bayesian regression based on material-specific feature vectors was used to predict the unknown coefficients along the phase-boundary lines as a function of the compositions. Similar approach was employed in the design of \ce{BaTiO3}-based electrocaloric ceramics.\cite{yuan2022optimizing}
\citet{balachandran2018experimental} trained the SVR model to predict the $T_C$ values, thus guiding the synthesis of high-temperature ferroelectric perovskites, more specifically, the materials system of the type Ba(MM$^\prime$)O$_3$-\ce{PbTiO3}, using feature vectors including information such as the tolerance factor, atomic valency, electronegativity, bond length and the Mendeleev number. 
\citet{yuan2022machine} employed the KRR model to predict the $d_{33}$ values to facilitate the search of \ce{BaTiO3}-based ceramics with large piezoelectric constants, where 95 different features that may affect the  $d_{33}$ values, including electronegativity, ionic radius, unit cell volume, ionic displacements and polarisation strengths, were used.\cite{yuan2018accelerated,yuan2019accelerated}
Finally, \citet{he2021machine} applied regression model to predict the permanent polarisations in ferroelectrics, using features including the atomic sizes and numbers, electrochemical factors, as well as the bond strength factors.

\textbf{\textit{Structure Classifications}}  Noncentrosymmetric crystals are proven to be difficult to discover experimentally, but are of great importance for applications such as nonlinear optics and piezoelectrics. \citet{song2020machine} applied the Random Forest classifier to predict the noncentrosymmetricity of inorganic solids, which helps to dicsover the mode relevant chemical features that can be correlated with the tendency of a material to be noncentrosymmetric. For this model, the author applied elemental features that were constructed using MATMINER\cite{ward2018matminer} and Magpie\cite{ward2016general} libraries. \citet{gleason2024random} developed an architecture of hierarchical random forest classifier to predict the crystal system, space group and lattice parameters from 2D electron diffraction patterns. Similar approach was adopted by \citet{corriero2023crystalmela} to learn the crystal structural information from the X-ray diffraction patterns. Finally, \citet{zhang2020phase} applied genetic algorithms to select the best ML model and the descriptors to distinguish between homogeneous solid solutions from two-phased mictures in high-entropy alloys, with descriptors that are composed of a wide range of elemental properties as well as the compositional information of the alloys. The fundamental intuition here is that the intermixing tendencies of the alloys can be determined by the compatibility in the chemical properties of the consituent elements. 

\textbf{\textit{Classification of Electronic Phases}} Based on a dataset that consists of thermally-driven metal-insulator transition (MIT) compounds, as well as metals and insulators with similar compositions and stoichiometries, \citet{georgescu2021database} applied the XGBoost\cite{chen2016xgboost} classifier model to distinguish compounds with these three distinguishable electronic phases, using features composes of the compositional, strutural, electronic and energetic properties of the compounds, with the feature importance further analysed with the SHAP method.\cite{lundberg2017unified} \citet{stanev2018machine} developed ML model to analyse the superconductor dadabase containing the $T_c$ values for 12,000 superconducting materials, in which the classifier model was trained to predict whether a material is a superconductor with $T_c$ above 10 K, using features that are constructed from stoichiometric information, elemental properties, and electronic structure properties. 


\paragraph{\textbf{Significance of Bayesian inference viewpoints for atomistic materials and molecular modelling}}

Before concluding this section, we would like to briefly discuss the significance of the Bayesian inference viewpoints for atomistic modelling, particularly for those that are backed by ML/AI models. Most importantly, we would like to stress that Bayesian inference is not just a mathematical trick for deriving the loss function, but deeply connected to the fundamental principles of statistical mechanics and thermodynamics, which are the cornerstones for understanding the behaviour of materials and molecules at the atomic level. 

First, we recite here the optimum model parameter for the ridge regression model derived in \cref{appendix:ridge_regression}:
\begin{equation}
    \label{eq:ridge_regression_optimum_parameter_bayesian_inference}
    \hat{\boldsymbol{\beta}}=\argmin_{\boldsymbol{\beta}}\left\{\mathbb{E}_{q}[\log p(\bm{y}|\boldsymbol{\beta})]-\mathcal{D}_{\mathrm{KL}}(q(\boldsymbol{\beta})|p(\boldsymbol{\beta}))\right\},
\end{equation}  
which essentially resembles the form of free energy in statistical mechanics. Here, the first term $\mathbb{E}_{q}[\log p(\bm{y}|\boldsymbol{\beta})]$ can be viewed as the negative of the internal energy, which quantifies how well the model predictions match the training data. Conceptually, the model is stable when the internal energy of the system is low, which is similar to the case when the ground truth can be well described by the model parameters. The second term $\mathcal{D}_{\mathrm{KL}}(q(\boldsymbol{\beta})|p(\boldsymbol{\beta}))$ can be viewed as the entropy of the system, which quantifies the uncertainty in the model parameters. Conceptually, the model is stable when the entropy of the system is high, which is similar to the case when the model parameters are not overly constrained by the training data and can capture a wide range of possible solutions to enhance the generalisability of the model. Therefore, training a regression model is a process that essentially balance the fidelity of the model to the training data and the complexity or the generalisability of the model, which is equivalent to the balance between the internal energy and the entropy in statistical mechanics.

Secondly, the Bayesian inference framework provides us with a systematic way to quantify the uncertainty in the model predictions, which, as we have seen for the development of on-the-fly PES sampling,\cite{jinnouchi2019phase} that maximises the trade-off between the accuracy and the computational cost in the computational process. 

Thirdly, the appropriate choice regularisation strategy in ML model training, as demonstrated in the LASSO regression [\cref{appendix:lasso_regression}], can be crucial for enhancing the interpretability of the model. This is an important aspect for atomistic modelling, in which enhancing our fundamental understanding of the underlying physical and chemical processes is often as important as making accurate predictions. It is particularly useful to reduce the model complexity thus enhance the model stabilities by focusing on a few key features that are most relevant to the target properties, as can be seen in the case of modelling the anharmonic vibrational behaviours of solids through the TDEP formalism\cite{hellman2013temperature} [\cref{sec:anharmonic_vibs}].

Finally, the feature vectors $\bm{\mathcal{X}}$ that are used as the input for ML models are  constructed based on the input structure $\bm{R}$ through some feature mapping functions $\bm{\mathcal{X}}=\bm{\mathcal{F}}_{\boldsymbol{\varphi}}(\bm{R})$, which are also parameterised by a set of hyperparameters $\boldsymbol{\varphi}$. In most ML models, these features are usually precomputed with a reasonable choice of $\boldsymbol{\varphi}$ due to the heaviness in computing these features, and the model training is only performed to optimise the model parameters $\boldsymbol{\theta}$. However, in principle, there is no barrier to further extend the loss function for the regression models to include the hyperparameters $\boldsymbol{\varphi}$, leading to the joint optimisation of $\mathcal{L}(\boldsymbol{\theta},\boldsymbol{\varphi})$ that simultaneously optimises the model parameters and the feature construction. Fundamentally, we can think of the application of deep learning models, such as deep NNs, that are applied to atomistic modelling as one that performs this process of joint optimisation.\cite{schutt2017schnet,gomez2018automatic,duval2023hitchhiker,hou2024unsupervised} 


